\documentclass[12pt]{article}
\newcommand{\level}{Grades 4--5}
\usepackage[letterpaper,top=0.45in,bottom=0.52in,left=0.5in,right=0.5in,
  includehead,includefoot,headheight=18pt,headsep=0.22in,footskip=0.34in]{geometry}
\usepackage[T1]{fontenc}
\usepackage[default]{sourcesanspro}
\usepackage{tikz}
\usetikzlibrary{arrows.meta}
\usepackage{fancyhdr}
\usepackage{array}
\pagestyle{fancy}
\fancyhf{}
\fancyhead[L]{Week 1 / Tiling / \level}
\fancyfoot[L]{Bellingham Math Circle / Week 1}
\fancyfoot[R]{\thepage}
\renewcommand{\headrulewidth}{0pt}
\renewcommand{\footrulewidth}{0pt}
\setlength{\parindent}{0pt}
\setlength{\parskip}{0pt}
\raggedright
\newcommand{\prob}[1]{\textbf{Problem #1:}}
\newcommand{\fig}[1]{\input{figs/#1.tex}}
\newcommand{\rarrow}{\tikz[baseline=-3pt]\draw[-{Stealth[length=9pt,width=8pt]},line width=1.3pt] (0,0)--(0.42in,0);}
\newcommand{\lrarrow}{\tikz[baseline=-3pt]\draw[{Stealth[length=9pt,width=8pt]}-{Stealth[length=9pt,width=8pt]},line width=1.3pt] (0,0)--(0.6in,0);}
\newcommand{\blank}[1][0.9in]{\rule{#1}{0.6pt}}
\newcommand{\green}{\raisebox{-1pt}{\fig{greenicon}}}
\newsavebox{\figbox}
\newcommand{\labfig}[2]{\sbox{\figbox}{\fig{#2}}\raisebox{\dimexpr\ht\figbox-1.6ex\relax}{\makebox[0.3in][l]{\textbf{#1}}}\usebox{\figbox}}

\newcommand{\chainrec}{\begin{tabular}[t]{@{}r@{\,}l@{}}left chain: & \blank[0.9in]\\[18pt] right chain: & \blank[0.9in]\end{tabular}}
\newcommand{\copyrow}[3]{\fig{#1}\foreach \i in {2,...,#2}{\hspace{#3}\fig{#1}}}
\begin{document}

In every problem, boards are covered with blue blocks only, with no gaps and no overlaps. Blank small drawings of boards are for recording.

\bigskip
\prob{1} Cover this board with blue blocks. Find every different way to do it, and draw each way in one of the small copies below. Explain how you know that you have found them all.

\vspace{0.3in}
\fig{h221}

\vspace{0.45in}
\centering
\labfig{A}{h221_s50}\hspace{0.35in}\labfig{B}{h221_s50}\hspace{0.35in}\labfig{C}{h221_s50}

\vspace{0.3in}
\labfig{D}{h221_s50}\hspace{0.35in}\labfig{E}{h221_s50}\hspace{0.35in}\labfig{F}{h221_s50}

\vspace{0.3in}
\labfig{G}{h221_s50}\hspace{0.35in}\labfig{H}{h221_s50}\hspace{0.35in}\hphantom{\labfig{H}{h221_s50}}

\raggedright
\newpage

\prob{2} When three blue blocks form a hexagon the shape of a yellow block, you may lift them and put them back the other way, as in the picture. This is one move. In the box, write the letters of your coverings from Problem 1, and draw a line between two letters whenever one move changes one covering into the other.

\vspace{0.3in}
\centering
\fig{moveA_g45}\hspace{0.25in}\raisebox{0.55in}{\lrarrow}\hspace{0.25in}\fig{moveB_g45}

\vspace{0.35in}
\framebox[\linewidth]{\rule{0pt}{4.6in}}

\raggedright
\newpage

\prob{3} Build the first covering of each pair on the big hexagon board below, then change it into the second covering using moves. Use as few moves as you can, and write how many you used.

\vspace{0.35in}
\fig{h222}

\vspace{0.5in}
\raisebox{1.35in}{\textbf{1}}\hspace{10pt}\fig{g45_p3_0_start}\hspace{0.2in}\raisebox{0.7in}{\rarrow}\hspace{0.2in}\fig{g45_p3_0_end}\hspace{0.5in}moves: \blank

\vspace{0.4in}
\raisebox{1.35in}{\textbf{2}}\hspace{10pt}\fig{g45_p3_1_start}\hspace{0.2in}\raisebox{0.7in}{\rarrow}\hspace{0.2in}\fig{g45_p3_1_end}\hspace{0.5in}moves: \blank

\newpage

\prob{4} Start with any covering of the big hexagon and make moves until you are back at the covering you started with. No move may undo the move just before it. Can this be done in exactly 4 moves? In 5? In 6? In 7? For each number that works, draw the coverings you pass through, in order. If you run out of copies, you can draw more on the back.

\vspace{0.3in}
4 moves: \blank[0.8in]\hfill 5 moves: \blank[0.8in]\hfill 6 moves: \blank[0.8in]\hfill 7 moves: \blank[0.8in]

\vspace{0.35in}
\centering
\copyrow{h222_s40}{4}{0.25in}

\vspace{0.2in}
\copyrow{h222_s40}{4}{0.25in}

\vspace{0.2in}
\copyrow{h222_s40}{4}{0.25in}

\vspace{0.2in}
\copyrow{h222_s40}{4}{0.25in}

\raggedright
\newpage

\prob{5} The bottom side of the big hexagon has two small edges. In each covering below, shade the block sitting on the left one of these edges. Then shade the block sitting on its top edge, and keep going until you reach the top side of the board. Do the same starting from the right edge. The blocks you shade from one edge form a chain. Each block in a chain leans left (L) or right (R), so a chain can be written as letters from bottom to top. Write both chains for each covering.

\vspace{0.2in}
\centering
\begin{tabular}{c@{\hspace{0.5in}}c}
\fig{blockL} & \fig{blockR}\\[4pt]
L & R
\end{tabular}

\vspace{0.3in}
{\setlength{\tabcolsep}{0pt}\begin{tabular}{@{}*{2}{>{\centering\arraybackslash}p{3.6in}}@{}}
\labfig{A}{g45_p5big_0} & \labfig{B}{g45_p5big_1}\\[14pt]
\chainrec & \chainrec\\
\rule{0pt}{0.5in} & \\
\labfig{C}{g45_p5big_2} & \labfig{D}{g45_p5big_3}\\[14pt]
\chainrec & \chainrec
\end{tabular}}

\raggedright
\newpage

\prob{6} For each pair of chains below, draw a covering of the big hexagon that has this left chain and this right chain, or explain why there is no such covering.

\vspace{0.3in}
\begin{tabular}{@{}p{3.6in}p{3.6in}@{}}
left chain LRLR, right chain RRLL & left chain RRLL, right chain LLRR\\[10pt]
\fig{h222_s45} & \fig{h222_s45}\\[0.35in]
left chain RLRL, right chain RLRL & left chain LRRL, right chain LRLR\\[10pt]
\fig{h222_s45} & \fig{h222_s45}
\end{tabular}

\vspace{0.5in}
\prob{7} When you make one move on a covering of the big hexagon, how can its chains change? Explain why a move always changes the chains in this way.

\newpage

\prob{8} For each pair of coverings in Problem 3, what is the fewest number of moves that changes the first covering into the second? Explain why it cannot be done in fewer moves.

\vspace{4.0in}
\prob{9} Start with a covering of the big hexagon, make some moves, and end at the covering you started with. Can the number of moves be odd? Explain.

\newpage

\prob{10} How many different ways can this board be covered with blue blocks? Explain how you know that you have counted every way. Small copies of the board for recording are on the next page.

\vspace{0.35in}
\fig{h133}

\newpage

\centering
\copyrow{h133_s30}{5}{0.3in}

\vspace{0.15in}
\copyrow{h133_s30}{5}{0.3in}

\vspace{0.15in}
\copyrow{h133_s30}{5}{0.3in}

\vspace{0.15in}
\copyrow{h133_s30}{5}{0.3in}

\vspace{0.15in}
\copyrow{h133_s30}{5}{0.3in}

\raggedright
\newpage

\prob{11} How many different ways can the big hexagon be covered with blue blocks? Explain how you know that you have counted every way. If you run out of copies, you can draw more on the back.

\vspace{0.3in}
\centering
\copyrow{h222_s36}{4}{0.3in}

\vspace{0.12in}
\copyrow{h222_s36}{4}{0.3in}

\vspace{0.12in}
\copyrow{h222_s36}{4}{0.3in}

\vspace{0.12in}
\copyrow{h222_s36}{4}{0.3in}

\vspace{0.12in}
\copyrow{h222_s36}{4}{0.3in}

\vspace{0.12in}
\copyrow{h222_s36}{4}{0.3in}

\end{document}
